Prospective students, parents, and visitors of a university campus face two
distinct information problems: understanding \emph{what} the institution
offers (admissions, programs, facilities, regulations) and understanding
\emph{where} things are on a campus they may never have physically visited.
Existing digital solutions typically address only one of these needs, and
conversational solutions built on large language models (LLMs) are prone to
fabricating institution-specific facts. This paper presents \system{}, a
full-stack platform that integrates a grounded retrieval-augmented
generation (RAG) assistant, a graph-structured 360$^{\circ}$ virtual tour,
and a database-backed A$^{\star}$ wayfinding engine into a single
application built on Next.js, FastAPI, PostgreSQL, ChromaDB, and a locally
hosted Llama~3.2 model served through Ollama. We describe, with reference to
the implemented source code, thirteen algorithmic components: a hybrid
dense--lexical retriever that fuses \code{all-MiniLM-L6-v2} cosine
similarity and BM25 with a $0.6/0.4$ weighting over min--max normalized
scores; a deterministic four-feature reranker; a two-term
intent/retrieval confidence estimator that gates generation; a
post-generation numeric-claim verification pass; a deterministic-first,
LSTM-fallback multi-agent router with five specialist agents; a
per-floor doubly linked list panorama scene graph with a tail-to-head
cross-floor stitching pass; a six-phase guided-tour traversal automaton;
an SVG minimap projection; and an A$^{\star}$ search whose heuristic
degrades gracefully from planar coordinates to a great-circle
approximation to zero (pure Dijkstra). We report implementation-scale
measurements: a $1{,}507$-chunk knowledge base derived from $128$ official
institutional PDFs and $42$ website pages, a $180$-node/$325$-edge campus
graph, $161$ panorama scenes across five floors, and a PyTorch LSTM intent
classifier trained on $168$ synthetic examples across $12$ classes
(training accuracy $1.00$, held-out validation accuracy $0.529$). We are
explicit that no formal retrieval-quality, answer-faithfulness, or
user-study evaluation has yet been performed; the paper therefore
contributes a system-level architecture, a detailed algorithmic
description, and a rigorous proposed evaluation protocol rather than
experimental performance claims.
